\documentclass[10pt]{article}

\usepackage{amsmath, amssymb, amsthm}
\usepackage{enumerate}
\usepackage{hyperref}
\usepackage{geometry}
\usepackage{newtxtext}
\usepackage{newtxmath}
\geometry{margin=1in}

\title{Banach Spaces - Lecture 1}
\author{}
\date{}

% ----------------------------------------------------------------
\begin{document}
\maketitle

% ================================================================

We present a proof of Pitt's theorem as a warm up.
The corresponding question is missing. But the proof is essentially all here.

\section*{Part 1. Warm Up}

\begin{itemize}
    

\item  From the question, let $\alpha = (\alpha_n)$ be a scalar sequence.

\item From the question, define $D_\alpha: \ell^p \rightarrow \ell^q$ by $D_\alpha(x) = (\alpha_n x_n)$ for $x = (x_n) \in \ell^p$.

\item We want to show that $D_\alpha$ is bounded if and only if $\alpha \in \ell^r$, where $\dfrac{1}{r} = \dfrac{1}{q} - \dfrac{1}{p}$.

\item Assume first that $\alpha \in \ell^r$ such that $\dfrac{1}{r} = \dfrac{1}{q} - \dfrac{1}{p}$.

\item By hypothesis on the $D_\alpha: \ell^p \rightarrow \ell^q$, we have that $||D_\alpha x||^q_q = \sum_{n=1}^\infty |\alpha_n|^q |x_n|^q$.

\item By manipulation of conjugate exponents we have $\dfrac{1}{r/q} + \dfrac{1}{p/q} = 1$. This was Cooly's explicit hint.

\item Continuing Cooly's explicit hint, applying H\"older's inequality with conjugate exponents $r/q$ and $p/q$, we have
\[||D_\alpha x||^q_q \leq \left( \sum_{n=1}^\infty (|\alpha_n|^q)^{r/q} \right)^{q/r} \left( \sum_{n=1}^\infty (|x_n|^q)^{p/q} \right)^{q/p} = ||\alpha||^q_r ||x||^q_p.\]

\item Taking the $q$-th root gives $||D_\alpha x||_q \leq ||\alpha||_r ||x||_p$, so $D_\alpha$ is bounded with operator norm at most $||\alpha||_r$. This gives the converse direction of the equivalence.

\item Assume now that $D_\alpha$ is bounded. We want to show that $\alpha \in \ell^r$.
This was a big pain for me since we cannot apriori assume $||\alpha||_r < \infty$.
For each $N$, define
\[
x^{(N)}_n=
\begin{cases}
|\alpha_n|^{r/p}\operatorname{sgn}(\alpha_n), & n\le N,\\
0,& n>N.
\end{cases}
\]

Let
\[
A_N=\sum_{n=1}^N |\alpha_n|^r.
\]

If $A_N>0$, put
\[
y^{(N)}=\frac{x^{(N)}}{A_N^{1/p}}.
\]

Then $\|y^{(N)}\|_p=1$, and a direct computation using
\[
\frac1r=\frac1q-\frac1p
\]
shows that
\[
\|D_\alpha y^{(N)}\|_q=A_N^{1/r}.
\]

Since $D_\alpha$ is bounded,
\[
A_N^{1/r}
=
\|D_\alpha y^{(N)}\|_q
\le
\|D_\alpha\|.
\]

Therefore
\[
\sum_{n=1}^N |\alpha_n|^r
\le
\|D_\alpha\|^r
\qquad (N\in\mathbb N).
\]

Letting $N\to\infty$ gives
\[
\sum_{n=1}^\infty |\alpha_n|^r<\infty,
\]
hence $\alpha\in\ell^r$.

The condition is required for H\"older's inequality to apply, so the forward direction is also true.

Since $(I - P_N)D_\alpha = D_{(0,\ldots,0, \alpha_{N+1}, \alpha_{N+2}, \ldots)}$ we get
\[
\|(I - P_N)D_\alpha\|_{\ell^p \to \ell^q} = \|(0, \dots, 0, \alpha_{N+1}, \alpha_{N+2}, \dots)\|_r = \|(\alpha_{N+1}, \alpha_{N+2}, \dots)\|_r
\]

This tells us that the tail of the sequence $\alpha$ must vanish as $N \to \infty$, so the theorem holds for diagonal operators.

\end{itemize}

\section*{Part 2.}

\begin{itemize}
    \item Let $K \subset \ell^q$ be a relatively compact set in $\ell^q$. Then, it is totally bounded since by definition, its closure is compact and subsets of compact sets are totally bounded. Note that $\ell^q$ is a complete metric space.
    \item By one characterisation of total boundedness, for any arbitrarily small $\epsilon > 0$, there exists a finite $\epsilon/2$-net for $K$. This means we can choose finitely many vectors $y_1, \dots, y_N \in K$ such that $K$ is contained in the union of finitely many balls of radius $\epsilon/2$ centered at $y_1, \dots, y_N$.
    \item Since each $y_i$ is in $\ell^q$ by assumption, we can choose $N$ such that $||(I - P_N)y_i||_q < \epsilon/2$.
    If $y \in K$, then pick $y_i$ such that $||y - y_i||_q < \epsilon/2$. 
    \item By the triangle inequality, $||(I-P_N)y||_q < \epsilon$. This proves the forward direction of the statement.
    \item For the converse, choose $N$ such that the $\mathrm{sup} ||(I - P_N)y_i||_q < \epsilon$, this allows us to reduce to the case of $P_N(K)$ in a finite dimensional subspace of $\ell^q$. 
    \item Using the above, consider $P_N(K)$ as a bounded subset of a finite dimensional subspace of $\mathbb{R}^N$. (our goal is to use the Heine-Borel theorem) Take the closure of this projection. 
    \item We use the following result without proof: the closure of the projection is closed and bounded since the closure of a bounded set remains bounded.
    \item By the Heine-Borel theorem, which we use without proof, the closure of $P_N(K)$ is compact and this is OK since we are working in $\mathbb{R}^N$.
    \item Using a definition of relative compactness, this implies $P_N(K)$ is relatively compact since its closure is compact, and we have $P_N(K)$ to be totally bounded.
    \item By one characterisation of total boundedness, for any arbitrarily small $\epsilon > 0$, there exists a finite $\epsilon/2$-net for $P_N(K)$. This means we can choose finitely many vectors $z_1, \dots, z_N \in P_N(K)$ such that $P_N(K)$ is contained in the union of finitely many balls of radius $\epsilon/2$ (dummy variable, lazy to change) centered at $z_1, \dots, z_N$.
    \item Since $P_N(K)$ is bounded, the $z_i$ form and $\epsilon/2$-net of $P_N(K)$, we can choose $z_i$ such that we have $||P_N y - z_i||_q < \epsilon/2$.
    \item By the triangle inequality, for any $y \in K$, we have $z_i$ such that $||y - z_i||_q < \epsilon$, by using the previous step and noting that $||(I - P_N)z_i||_q < \epsilon/2$ as hypothesis.
    \item Since the $z_i$ forms a finite $\epsilon$-net, $K$ is totally bounded.
    \item Since $K$ is totally bounded, it is relatively compact. This proves the converse direction of the statement.
    \item We compare this to the original problem. Therefore, we simply need to show the image of $T$ of the unit ball is relatively compact in $\ell^q$.
\end{itemize}

\section*{Part 3.}

\begin{itemize}
    \item From the question, let $1 \leq r < \infty$. Suppose $u_n \to 0$ weakly in $\ell^r$. Fix $a \in \ell^r$, and consider the case for $a$ with finite support.
    \item By hypothesis of finite support, suppose $a_k = 0$ for all $k > M$. We split the norm into two parts:
    \[
       ||a+u_n||^r_r = \sum_{k=1}^M |a_k + u_{n,k}|^r + \sum_{k=M+1}^\infty |u_{n,k}|^r.
    \]
    \item By hypothesis of the question, we have $u_n \to 0$ weakly in $\ell^r$. Then we have $u_{n,k}$ having coordinate-wise convergence to $0$ for each $k$.
    \item By examining the first part of the split of the norm, this reduces the part within the support to $||a||^r_r$, which is what we want later.
    \item For the part outside of the support, rewrite it in terms of total norm of $u_n$ (we want to use the hypothesis of weak convergence later) so that:
    \[
       \sum_{k=M+1}^\infty |u_{n,k}|^r = ||u_n||^r_r - \sum_{k=1}^M |u_{n,k}|^r.
    \]
    \item For the norm split into two parts, take limit supremum on both sides, use the weak convergence of $u_n \to 0$ weakly in $\ell^r$ to kill off the term involving $u_n$. Then we have $u_{n,k}$ having coordinate wise convergence to $0$:
    \[
       \limsup_{n \to \infty} ||a+u_n||^r_r
       = ||a||^r_r + \limsup_{n \to \infty} ||u_n||^r_r.
    \]
    which is our desired result for the case of finite support.
    \item Now take a general $a \in \ell^r$. We can approximate $a$ by a sequence of finite support vectors $a^m$ such that $||a - a^m||_r \to 0$ as $m \to \infty$. Apply the finite support result to $a^m$ and pass to the limit as $m \to \infty$ to get the desired result for general $a$.
\end{itemize}

\section*{Part 4.}

\begin{itemize}
    \item Using the Cooly hint, suppose for the sake of contradiction $||T h_n||_q \not\to 0$ as $n \to \infty$. Without loss of generality, for some subsequence (it inherits weak convergence), we can assume $||T h_n||_q > L$ for some $L > 0$ and all $n$.
    \item Using the hypothesis that $(h_n)$ is bounded, there exists an $M > 0$ such that $||h_n||_p < M$ for all $n$.
    \item Since $||T|| = 1$ by normalisation, pick a unit vector $x$ in $\ell^p$ such that $T$ nearly attains its norm i.e. $||T x||_q > 1 - \epsilon$. Using the hint, perturb it with small $h_n$, consider the case $||T(x+th_n)||_q = ||Tx + t T h_n||$.
    \item The weakly null sequence $(h_n)$ satisfies the hypotheses, work in the domain space $\ell^p$:
    \[
    \limsup_{n \to \infty} \|x + t h_n\|_p^p = \|x\|_p^p + t^p \limsup_{n \to \infty} \|h_n\|_p^p \leq \|x\|_p^p + t^p M^p
    \]
    \item Using the hypothesis that $T$ is a bounded linear operator, and $1 \leq q$, and we use the fact that bounded linear operators map weakly null sequences to weakly null sequences, now we can work in the image space $\ell^q$,and apply part 3 again. Apply the first Cooly hint for contradiction to make an assumption on $L$, and apply the other Cooly hint and have $T$ nearly obtain its norm, this gives the bound:
    \[
    \limsup_{n \to \infty} \|T x + t T h_n\|_q^q = \|Tx\|_q^q + t^q \limsup_{n \to \infty} \|Th_n\|_q^q \geq 
    (1 - \epsilon)^q + t^q L^q
    \] 
    On the other hand,
    \[
    \|Tx+tTh_n\|_q
    =
    \|T(x+th_n)\|_q
    \le
    \|T\|\,\|x+th_n\|_p
    =
    \|x+th_n\|_p,
    \]
    since $\|T\|=1$.

    Taking $\limsup$ and using the estimate obtained in the domain space gives
    \[
    \limsup_{n\to\infty}\|Tx+tTh_n\|_q^q
    \le
    \left(
    1+t^pM^p
    \right)^{q/p}.
    \]

    Combining the two inequalities yields
    \[
    (1-\varepsilon)^q+t^qL^q
    \le
    (1+t^pM^p)^{q/p}.
    \]
    \item Use the normalisation $||T|| = 1$ and our definition of $M$, we have by cheating and looking at what we want for part 5:
    \[
        (1 - \epsilon)^q + t^q L^q \leq (1 + t^p M^p)^{q/p}
    \]
    \item Isolate $L^q$, we get:
    \[
    L^q
    \le
        \frac{(1+t^pM^p)^{q/p}-(1-\varepsilon)^q}{t^q}.
     \]
    \item Looking at Part 5, we just need to show $L$ = 0, get contradiction that $L > 0$. Assume part 5 holds, we will have proven Part 4.
\end{itemize}

\newpage

\section*{Part 5}

\begin{itemize}
    \item Assume $0 < q < p < \infty$ and $L \geq 0$. We have \[
    L^q
    \le
        \frac{(1+t^pM^p)^{q/p}-(1-\varepsilon)^q}{t^q}.
     \]
    we want to show $L = 0$.
    \item Fix $t > 0$ without loss of generality, let $\epsilon \rightarrow 0^+$
    \[
    L^q
    \le
        \frac{(1+t^pM^p)^{q/p}-1}{t^q}.
    \]
    \item For small $t$, Taylor's expansion gives $(1+t^pM^p)^{q/p} = 1 + \dfrac{q}{p} M^p t^p + o(t^p)$
    Reduce to $L^q \leq C t^{p-q} + o(t^{p-q})$.
    \item By hypothesis $p - q > 0$, let $t \rightarrow 0$ gives $L^q = 0$ so $L = 0$.
    \item Part 4 is proven as well.
\end{itemize}

\section*{Part 6}

\begin{itemize}
    \item From the main statement $1 < q < p < \infty$, $\ell^p$ is reflexive. Using a characterisation of the Banach-Alaoglu theorem, a corollary is that every bounded set in a reflexive Banach space is relatively weakly compact. We just need that every bounded sequence has a weakly convergent subsequence.
    \item By part 2, it suffices to show that the image of operator $T$ of the unit ball (denote $K$) is relatively compact in $\ell^q$.
    \item Consider $(y_n)$ in $K$, say $y_n = T x_n$ with $||x_n||_p < 1$. By the reflexivity of $\ell^p$ there is a subsequence $x_{n_k}$ weakly converging to $x$ in $\ell^p$. 
    \item Set $h_k := x_{n_k} - x$. Then $h_k$ tends to $0$ weakly in $\ell^p$ by linearity and $(h_k)$ is bounded. We can apply parts $4$ and $5$ (normalise if needed):
    \[
    \|Tx_{n_k}-Tx\|_q
    =
    \|Th_k\|_q
    \to 0.
    \]
    So the subsequence $y_{n_k} = T x_{n_k}$ converges to $Tx$ in $\ell^q$, every sequence in $K$ has a convergent subsequence, $K$ is sequentially relatively compact, and in a metric space, therefore relatively compact.
    \item By part 2, the theorem is proved.
\end{itemize}

\end{document}